\chapter{Tổng quan đề tài}

\section{Bối cảnh}

Trong thực tế, nhiều tổ chức làm việc với số lượng lớn tài liệu nội bộ như hợp đồng, biểu mẫu phê duyệt, tài liệu kỹ thuật, hồ sơ nhân sự, tài liệu kiểm toán hoặc các file phục vụ quá trình vận hành. Các tài liệu này thường không chỉ được lưu một lần rồi để đó, mà sẽ đi qua nhiều bước xử lý khác nhau như tạo mới, cập nhật phiên bản, gửi duyệt, phê duyệt, xuất bản, truy cập, tải xuống hoặc lưu trữ. Nếu hệ thống chỉ lưu file mà không quản lý trạng thái và lịch sử xử lý, việc kiểm tra lại tài liệu sau này sẽ gặp nhiều khó khăn.

Một vấn đề thường gặp là khi cần xem lại một tài liệu, người quản trị hoặc người phụ trách không chỉ muốn biết file đang nằm ở đâu, mà còn muốn biết tài liệu đó đã được ai tạo, ai cập nhật, đã qua phiên bản nào, ai phê duyệt, ai từng truy cập và hiện tại đang ở trạng thái nào. Với tài liệu nhạy cảm, các thông tin này càng quan trọng vì nó liên quan đến trách nhiệm xử lý, kiểm soát truy cập và khả năng chứng minh tài liệu đã được quản lý đúng quy trình.

Từ nhu cầu đó, nhóm xây dựng DocVault như một hệ thống quản lý tài liệu bảo mật. Trọng tâm của phần ứng dụng là quản lý tài liệu theo vòng đời và ghi nhận lại các hoạt động quan trọng bằng audit log. Mỗi thao tác như tạo tài liệu, upload version, gửi duyệt, phê duyệt, preview, download, retention hoặc compliance đều cần để lại dấu vết để khi cần có thể kiểm tra lại. Nhờ đó, tài liệu trong hệ thống không chỉ là một file lưu trữ, mà là một đối tượng có metadata, trạng thái, quyền truy cập, lịch sử và bằng chứng xử lý.

Bên cạnh phần ứng dụng, đề tài cũng xây dựng phần DevSecOps. Lý do là một hệ thống microservices gồm nhiều service, nhiều database, object storage, Kubernetes manifest và các cơ chế bảo mật sẽ rất khó triển khai ổn định nếu làm hoàn toàn thủ công. Khi có thay đổi mã nguồn hoặc cấu hình, nhóm cần một quy trình có thể tự động kiểm thử, quét bảo mật, build image, đẩy artifact, cập nhật GitOps và kiểm tra lại trạng thái sau triển khai. Đây là lý do DevSecOps được xem là phần chính thứ hai của đề tài.

Như vậy, DocVault được xây dựng theo hai hướng song song. Phần ứng dụng trả lời câu hỏi hệ thống quản lý tài liệu bảo mật được thiết kế và vận hành như thế nào. Phần DevSecOps trả lời câu hỏi hệ thống đó được kiểm thử, đóng gói, triển khai và kiểm chứng ra sao. Hai phần này bổ sung cho nhau để tạo thành một hệ thống vừa có chức năng nghiệp vụ, vừa có quy trình phát hành và vận hành có kiểm soát.

\section{Vấn đề cần giải quyết}

Đề tài đặt ra hai nhóm vấn đề chính. Nhóm thứ nhất liên quan đến việc xây dựng ứng dụng quản lý tài liệu bảo mật. Nhóm thứ hai liên quan đến việc xây dựng quy trình DevSecOps để kiểm thử, triển khai và kiểm chứng hệ thống.

Đối với phần ứng dụng, vấn đề chính là làm sao quản lý được toàn bộ vòng đời của một tài liệu. Một tài liệu có thể được tạo mới, upload phiên bản, cập nhật metadata, gửi duyệt, phê duyệt, từ chối, xuất bản, truy cập, tải xuống, lưu trữ hoặc đưa vào retention. Nếu chỉ lưu file và trạng thái hiện tại, hệ thống sẽ khó trả lời các câu hỏi như tài liệu này đã được xử lý qua những bước nào, ai đã thao tác, thao tác đó có hợp lệ không và có bằng chứng nào để kiểm tra lại.

Vì vậy, hệ thống cần quản lý đồng thời nhiều thông tin liên quan đến tài liệu, bao gồm metadata, version file, checksum, trạng thái workflow, quyền truy cập, mức phân loại bảo mật, lịch sử phê duyệt và audit log. Các thông tin này cần được liên kết với nhau để tạo thành một bức tranh đầy đủ về quá trình xử lý tài liệu. Khi cần kiểm tra lại, hệ thống phải có khả năng tổng hợp các dữ liệu này thành evidence để phục vụ việc đối chiếu.

Đối với phần DevSecOps, vấn đề đặt ra là làm sao đưa một hệ thống nhiều service vào một quy trình triển khai có thể lặp lại và kiểm soát được. Nếu build image, cập nhật manifest, triển khai Kubernetes và kiểm tra trạng thái đều làm thủ công, quá trình phát hành sẽ dễ sai sót và khó truy vết. Do đó, hệ thống cần pipeline để chạy kiểm thử, quét bảo mật, build image, phát hành artifact, cập nhật GitOps, đồng bộ qua Argo CD và kiểm tra lại sau triển khai.

Cụ thể, đề tài tập trung vào các yêu cầu sau:

\begin{itemize}
  \item Quản lý vòng đời tài liệu với các trạng thái như \code{DRAFT}, \code{PENDING}, \code{PUBLISHED} và \code{ARCHIVED}.
  \item Quản lý metadata, version file, checksum, workflow history và trạng thái hiện tại của tài liệu.
  \item Kiểm soát quyền truy cập dựa trên RBAC, ACL, classification, trạng thái tài liệu và quyền sở hữu.
  \item Ghi nhận audit log cho các thao tác quan trọng như tạo tài liệu, upload version, submit, approve, reject, preview, download, retention và compliance.
  \item Cung cấp evidence packet để tổng hợp metadata, version checksum, workflow history, retention evidence và audit event liên quan đến tài liệu.
  \item Kiểm soát preview/download thông qua backend policy và grant token ngắn hạn.
  \item Đảm bảo Compliance Officer chỉ xem metadata, audit và evidence phục vụ kiểm tra, không truy cập trực tiếp nội dung file.
  \item Xây dựng pipeline DevSecOps để kiểm thử, quét bảo mật, build image, đẩy artifact lên Harbor, cập nhật GitOps, đồng bộ qua Argo CD, chạy smoke test, ZAP và thu thập bằng chứng vận hành.
\end{itemize}

\section{Mục tiêu đề tài}

Mục tiêu tổng quát của đề tài là xây dựng và triển khai hệ thống DocVault gồm hai phần chính: phần ứng dụng quản lý tài liệu bảo mật và phần DevSecOps cho quá trình kiểm thử, triển khai, vận hành hệ thống.

Đối với phần ứng dụng, nhóm hướng đến việc xây dựng một hệ thống có thể quản lý tài liệu theo vòng đời rõ ràng, có kiểm soát truy cập và có audit/evidence để kiểm tra lại quá trình xử lý. Đối với phần DevSecOps, nhóm hướng đến việc tự động hóa các bước kiểm thử, quét bảo mật, build, phát hành artifact, cập nhật GitOps, triển khai lên Kubernetes và kiểm tra sau triển khai.

Các mục tiêu cụ thể gồm:

\begin{enumerate}
  \item Thiết kế kiến trúc tổng thể cho hệ thống DocVault, gồm frontend, API Gateway, các backend service, các thành phần lưu trữ dữ liệu và hạ tầng triển khai.
  \item Xây dựng các chức năng quản lý vòng đời tài liệu như tạo tài liệu, cập nhật metadata, upload version file, gửi duyệt, phê duyệt, từ chối, xuất bản, lưu trữ và retention.
  \item Thiết kế mô hình dữ liệu cho documents, document versions, ACL, workflow history, audit events và các thông tin evidence liên quan.
  \item Xây dựng cơ chế kiểm soát truy cập dựa trên RBAC, ACL, trạng thái tài liệu, classification, quyền sở hữu và quy tắc deny-overrides.
  \item Ghi nhận audit log cho các thao tác quan trọng trong vòng đời tài liệu và áp dụng hash-chain để hỗ trợ kiểm tra tính toàn vẹn của audit event.
  \item Xây dựng evidence packet để tổng hợp metadata, checksum phiên bản, workflow history, retention evidence và audit event, giúp kiểm tra lại quá trình quản lý tài liệu.
  \item Tích hợp các cơ chế bảo mật như grant token ngắn hạn cho preview/download, Compliance Officer metadata-only boundary, malware scanning và DLP.
  \item Xây dựng hạ tầng triển khai cloud-native trên AWS EKS với Terraform, Helm/Kustomize, S3/KMS, IRSA, External Secrets, Harbor và các thành phần quan sát vận hành.
  \item Xây dựng pipeline DevSecOps với Jenkins và các công cụ kiểm thử/quét bảo mật như SonarQube, Dependency Check, Trivy, Checkov, Terraform Validate, Kyverno, Cosign, Argo CD, smoke test và OWASP ZAP.
  \item Thu thập bằng chứng cho cả phần ứng dụng và phần DevSecOps, bao gồm E2E test, audit/evidence, retention, scan report, image digest, SBOM, signing status, GitOps revision, Argo CD health, smoke test và ZAP report.
\end{enumerate}

\section{Phạm vi thực hiện}

Phạm vi đề tài gồm hai phần có quan hệ chặt chẽ. Phần ứng dụng bao gồm giao diện web, các dịch vụ backend, luồng nghiệp vụ tài liệu, kiểm soát truy cập, audit, DLP, quét mã độc, retention và evidence center. Phần DevSecOps bao gồm hạ tầng AWS EKS được quản lý bằng Terraform, quy trình CI/CD trên Jenkins, kho ảnh Harbor, GitOps với Argo CD, các cổng kiểm soát bảo mật và cơ chế giám sát sau triển khai.

Tại thời điểm hoàn thành báo cáo, pipeline đã tích hợp SAST, SCA, quét secret, Trivy, Checkov, Terraform Validate, Kyverno CLI, build và quét container image, đẩy artifact lên Harbor, cập nhật GitOps, kiểm tra trạng thái Argo CD, smoke test và DAST bằng OWASP ZAP. Hạ tầng triển khai đã sử dụng Amazon S3 với SSE-KMS, IAM Role for Service Account (IRSA), AWS Secrets Manager kết hợp External Secrets Operator, backup cơ sở dữ liệu, Prometheus, Grafana và Loki. Harbor đã lưu thông tin SBOM và trạng thái ký của artifact. Các nội dung còn lại tập trung vào vận hành ở quy mô production, thay vì bổ sung các thành phần nền tảng của quy trình.

\begin{table}[H]
\centering
\caption{Phạm vi và trạng thái thực hiện của đề tài}
\label{tab:scope-status}
\begin{tabularx}{\textwidth}{p{0.28\textwidth}Yp{0.22\textwidth}}
\toprule
\textbf{Nhóm nội dung} & \textbf{Mô tả} & \textbf{Trạng thái} \\
\midrule
Ứng dụng web & Frontend Next.js, các màn hình dashboard, documents, approvals, audit, security, evidence, retention. & \statusdone \\
Backend microservices & Gateway, metadata-service, document-service, workflow-service, audit-service, notification-service. & \statusdone \\
Bảo mật runtime & RBAC, ACL, grant token, compliance boundary, DLP, malware scanning, audit hash-chain. & \statusdone \\
Local infrastructure & Postgres, MongoDB, MinIO, Keycloak qua Docker Compose. & \statusdone \\
DevSecOps pipeline & Jenkins pipeline với các stage kiểm thử, quét bảo mật, build, Harbor, GitOps, Argo CD, smoke test và DAST. & \statusdone \\
Hạ tầng cloud-native & Terraform, EKS, S3/KMS, IRSA, External Secrets, backup, Harbor và observability. & \statusdone \\
Gia cố vận hành production & Remote Terraform state, private node, admission enforcement, provenance chuẩn hóa, NetworkPolicy giới hạn egress, log persistence và diễn tập khôi phục. & \statusprogress \\
\bottomrule
\end{tabularx}
\end{table}

\section{Đóng góp của đề tài}

Đề tài có bốn đóng góp. Thứ nhất là một hệ thống quản lý tài liệu chạy được, bao phủ luồng từ khởi tạo, tải lên, phê duyệt đến lưu trữ. Thứ hai là mô hình kiểm soát truy cập nhiều lớp, kết hợp vai trò, ACL, mức phân loại, trạng thái và quyền sở hữu thay vì chỉ dựa vào một trong số đó. Thứ ba là cách đưa yêu cầu compliance vào ngay trong thiết kế: audit hash-chain để phát hiện sửa đổi, retention evidence và evidence packet metadata-only. Thứ tư là quy trình DevSecOps được nhóm đầu tư triển khai khá đầy đủ, có khả năng truy vết từ mã nguồn qua container artifact và GitOps revision đến workload trên EKS, với các cổng kiểm soát chất lượng, bảo mật chuỗi cung ứng, quản lý bí mật và giám sát vận hành gắn vào từng chặng.

\section{Phân công công việc}

Đề tài do hai thành viên cùng thực hiện theo hình thức đồng phụ trách: cả hai cùng tham gia khảo sát, thiết kế và kiểm thử toàn hệ thống, đồng thời mỗi người nhận vai trò chính ở một số hạng mục để bảo đảm tiến độ. Bảng~\ref{tab:work-assignment} tóm tắt phân công theo vai trò chính và vai trò phối hợp.

\begin{table}[H]
\centering
\small
\caption{Phân công công việc giữa các thành viên}
\label{tab:work-assignment}
\begin{tabularx}{\textwidth}{p{0.34\textwidth}p{0.26\textwidth}Y}
\toprule
\textbf{Hạng mục} & \textbf{Phụ trách chính} & \textbf{Phối hợp} \\
\midrule
Thiết kế kiến trúc tổng thể & Cả hai & --- \\
Frontend Next.js và permission-aware UI & Huỳnh Lê Đại Thắng & Nguyễn Trường Duy \\
Backend microservices và bảo mật runtime & Huỳnh Lê Đại Thắng & Nguyễn Trường Duy \\
Mô hình phân quyền (RBAC, ACL, classification) & Huỳnh Lê Đại Thắng & Nguyễn Trường Duy \\
Audit hash-chain, DLP và evidence & Cả hai & --- \\
Pipeline DevSecOps (Jenkins, các security gate) & Nguyễn Trường Duy & Huỳnh Lê Đại Thắng \\
Hạ tầng cloud-native (Terraform, EKS, GitOps) & Nguyễn Trường Duy & Huỳnh Lê Đại Thắng \\
Quản lý secret, registry và observability & Nguyễn Trường Duy & Huỳnh Lê Đại Thắng \\
Kiểm thử, thu thập bằng chứng và viết báo cáo & Cả hai & --- \\
\bottomrule
\end{tabularx}
\end{table}

\section{Cấu trúc báo cáo}

Báo cáo gồm tám chương chính. Chương 1 trình bày bối cảnh, mục tiêu và phạm vi đề tài. Chương 2 trình bày cơ sở lý thuyết và công nghệ nền tảng. Chương 3 phân tích yêu cầu và thiết kế tổng thể. Chương 4 mô tả thiết kế dữ liệu, API và chính sách bảo mật. Chương 5 trình bày quá trình triển khai ứng dụng DocVault. Chương 6 trình bày quy trình DevSecOps và hạ tầng triển khai. Chương 7 mô tả kiểm thử, đánh giá và bằng chứng. Chương 8 tổng kết kết quả đạt được, hạn chế và hướng phát triển.
